\subsection{若干精确稳态解}
\label{sec:chap1-exact-steady-solutions}

在后续各章精确分析方程之前, 本节给出若干存在显式公式的稳态例子. 我们求解齐次不可压缩 Navier--Stokes 方程\eqref{eq:chap1-dimensionless-navier-stokes}的稳态形式
\begin{equation}
\label{eq:chap1-steady-dirichlet-system}
\left\{\begin{aligned}
&-\frac1{\operatorname{Re}}\Delta v+(v\cdot\nabla)v+\nabla p=g&&\text{在 $\Omega$ 中},\\
&\operatorname{div}v=0&&\text{在 $\Omega$ 中},\\
&v=v_b&&\text{在 $\partial\Omega$ 上}.
\end{aligned}\right.
\end{equation}
其中 $g$ 是浮力, $v_b$ 是 Dirichlet 边界数据. 若 $\nabla\pi_0=g$, 令 $\pi=p-\pi_0$ 即可化为零源项情形. 显式解要求区域具有简单几何. 当 Reynolds 数较大时, 这些解析解未必就是实际观察到的解, 还需研究稳定性; 本书不讨论这一重要问题, 参见\MBUnresolvedReference{bib:stability-references}{[40, 53, 127, 43]}.

\subsubsection{管内 Poiseuille 流}
\label{subsec:chap1-poiseuille-flow}

考虑轴向为 $e_3$ 且截面为椭圆的无限长管. 截面 $E$ 定义为
\begin{equation}
\label{eq:chap1-elliptic-pipe-section}
  \left(\frac{x_1}{R_1}\right)^2+\left(\frac{x_2}{R_2}\right)^2\leq1.
\end{equation}
由对称性设 $v=(u_1(x_1,x_2),u_2(x_1,x_2),u_3(x_1,x_2))$. 令 $\widetilde v=(u_1,u_2)$, 则 $\widetilde v$ 满足二维稳态 Navier--Stokes 方程和边界条件 $\widetilde v=0$. 用 $\widetilde v$ 乘方程并在 $E$ 上积分, 分部积分及无散条件给出
\[
  \frac1{\operatorname{Re}}\int_E(|\nabla u_1|^2+|\nabla u_2|^2)\,\mathrm{d}x_1\mathrm{d}x_2=0.
\]
故 $u_1=u_2=0$, $\pi$ 与 $x_1,x_2$ 无关, 且 $v=(0,0,u(x_1,x_2))$. 此时 $(v\cdot\nabla)v=0$. 记常数 $\partial\pi/\partial x_3=p_1$, 则
\[
  -\frac1{\operatorname{Re}}\Delta u+p_1=0\quad\text{在 $E$ 中},
  \qquad u=0\quad\text{在 $\partial E$ 上}.
\]
唯一解为
\[
u=\frac{-\operatorname{Re}p_1}{R_1^{-2}+R_2^{-2}}
\left(1-\frac{x_1^2}{R_1^2}-\frac{x_2^2}{R_2^2}\right),
\qquad \pi=p_1x_3.
\]
所以压力梯度驱动流动. 实际压力为 $p=\pi+\pi_0$, 其中 $\pi_0$ 是重力势.

\subsubsection{平面剪切流}
\label{subsec:chap1-planar-shear}

考虑两块水平平板之间的二维 Newtonian 黏性流体. 上板以速度 $U\geq0$ 运动, 下板以速度 $-V\leq0$ 运动, $\Omega=\mathbb{R}\times(-h,h)$. 几何和无散条件给出 $v=(u(x_2),0)$. 惯性项自动为零, 且
\[
  -\frac1{\operatorname{Re}}\Delta u+\frac{\partial\pi}{\partial x_1}=0,
  \qquad\frac{\partial\pi}{\partial x_2}=0.
\]
记常数 $\partial\pi/\partial x_1=p_1$, 并用 $u(h)=U$, $u(-h)=-V$, 得
\[
u=\frac{\operatorname{Re}p_1}{2}(h^2-x_2^2)+\frac{U+V}{2h}x_2+\frac{U-V}{2},
\qquad \pi=p_1x_1.
\]
这里压力梯度和边界条件共同驱动流动, 实际压力仍为 $p=\pi+\pi_0$.

\subsubsection{两圆柱间的 Couette 流}
\label{subsec:chap1-couette-flow}

考虑由两个轴向为 $e_3$ 的同轴无限长圆柱限定的区域. 外圆柱固定, 内圆柱以与 $e_3$ 共线的旋转向量 $\Omega$ 匀速旋转. 仿照刚体旋转速度 $v=\Omega\wedge r$, 设
\[
  v=f(r)\Omega\wedge e_r.
\]
该向量场无散. 利用向量恒等式和柱坐标公式,
\begin{equation}
\label{eq:chap1-couette-curl}
  \operatorname{curl}v=\frac1r\frac{\partial}{\partial r}(rf(r))\Omega.
\end{equation}
从而 $v\wedge\operatorname{curl}v$ 是梯度, 惯性项可写为 $(v\cdot\nabla)v=\nabla\widetilde\pi(r)$, 其中
\[
  \widetilde\pi'(r)=-\frac{|\Omega|^2}{r}f^2(r).
\]
稳态方程化为
\[
  -\frac1{\operatorname{Re}}\Delta(\Omega\wedge f(r)e_r)+\nabla(\pi+\widetilde\pi)=0.
\]
取 $\Delta(\Omega\wedge f(r)e_r)=0$ 和 $\pi=-\widetilde\pi$. 若 $F'=f$, 则 $\Delta F=a$, 因而
\[
  F(r)=\frac a4r^2+b\ln r+c,\qquad f(r)=\frac a2r+\frac br.
\]
令内外圆柱半径分别为 $R_1,R_2$, 边界条件 $f(R_1)=R_1$, $f(R_2)=0$ 给出
\[
  f(r)=\frac{R_1^2R_2^2}{R_2^2-R_1^2}\left(\frac1r-\frac r{R_2^2}\right),
  \qquad v=f(r)|\Omega|e_\theta.
\]
函数 $f$ 递减且为凸函数. 压力为
\[
\pi(r)=|\Omega|^2\left(\frac{R_1^2R_2^2}{R_2^2-R_1^2}\right)^2
\left(\frac{r^2}{2R_2^4}-\frac1{2r^2}-\frac2{R_2^2}\ln r\right).
\]
它随 $r$ 增加, 因而外圆柱附近压力高于内圆柱附近压力. 实际压力为 $p=\pi+\pi_0$. 与前两个例子不同, 此解析解与 Reynolds 数无关: 惯性项恰被压力项抵消, 黏性项恒为零.
